% !TeX root = ../papers/nc-p3.tex
\section{十八次の核グラフ切断と滑らかな三次核曲面の排除}
\label{sec:cofactor-eighteen-sectional-graph}
\begin{theorem}\label{thm:cofactor-sectional-graph-genus}
整で一般被約な階数二曲線の核像 $K$ が曲面であるとする。一般入力平面での正規化を $B$、
入力平面と核空間の積への像を $G$（被約像）、核空間での像を $J$ と書く。
$G\to J$ は有限双有理であり、$p_a(G)=21-d$ である。
特に $d=18$ なら $J$ は非退化な整六次曲線で $p_a(J)=3,4$、
$p_a(J)=3$ なら $G=J$ である。滑らかな核曲面に双有理で $p_a(J)=3$ なら
逆四空間の基点イデアルは全て局所完全交叉である。
\end{theorem}
\begin{proof}
核グラフ曲面から $K$ への射の一般繊維は有限である。相異なる二入力点が同じ一般核点を
持つ対は高々二次元族で、一対を含む入力平面は一次元族である。
従って一般平面上にそのような対は高々有限個である（特別な核点を避けた開集合で十分）。
よって $G\to J$ は非定数の双有理な射影曲線射である。入力平面八次 $E_H$ への射
$h:G\to E_H$ も有限双有理である。階数三で繊維は線形核一点に埋まり、
階数二の横断点では長さ二以下である。後者は一次元Cohen--Macaulay加群
$h_*\mathcal O_G$ の最小生成数が一般パラメータで割った長さを超えず、その長さが
階数一と $E_H$ の重複度二から二になることによる。各繊維は核空間に閉埋込みであり、
長さ二以下なら一次式で生成される。従って平面上の評価
$M_H\to h_*\mathcal O_G(0,1)$ は全射、純性と一般同型から同型。
Euler数より $4=24-d+1-p_a(G)$ を得る。$d=18$ のとき $\deg K\le4$ なので
六次の $J$ は平面に入らず、Castelnuovo上界は四である。

滑らかな双有理の場合、一般平面はある逆切断 $a$ に対応し、$J=(a=0)$ である。
$G=J$ なら入力の四切断の比の射が $J$ に正則に延長される。
各基点で $I_Z\mathcal O_J$ は従って主イデアルである。実際射の像で単元となる射影座標を選べば
全座標はその切断の正則な関数倍となる。その切断を $b$ として持ち上げれば
$S_0$ の局所環で $I_Z=(a,b)$。共通曲線はなく正則列である。長さの和は
一般二切断の交点数から一般入力直線の逆像長八を引いた $J^2-8$ となる。
\end{proof}

\begin{theorem}\label{thm:cofactor-eighteen-smooth-cubic}
$d=18$ で核像は滑らかな三次曲面にならない。また有理二重点のみの正規四次曲面
（滑らかな場合を含む）も核像にならない。
\end{theorem}
\begin{proof}
三次または四次なら定理\ref{thm:cofactor-sectional-kernel}から核写像は双有理。
四次の場合、最小解消で $K_{S_0}=0$。固定成分を除いた逆系 $A$ は
$AL=6$, $8\le A^2\le9$ を満たし、随伴で偶数なので $A^2=8$。基点なしで
一般切断の種数は五となり上界三に反する。

滑らかな三次曲面 $S_0$ では $-K_{S_0}=L$, $AL=6$ で
$p_a(J)=1+(A^2-6)/2$。従って $A^2=10,12$。十なら単純二基点（無限近接も許す）。
十二ならHodge等号から $A=2L$（有理曲面のPicard群に数値的自明な非零元はない）。
四単純基点では一般切断の正規化種数が四となるので、一つの二重基点だけである。
いずれも全基点の解消 $S$ は滑らかで $H^2=8$, $HK_S=-4$, $HL=6$。
$\chi(S,L+tH)=4(t+1)^2$ で、定理\ref{thm:cofactor-exclude-nineteen}と同様、
入力射が有限な開集合で $M=h_*L$ が成立する。

単純二基点の場合 $P=A-L$ と置くと $P^2=1$, $PL=3$, $R=P-L$ は
$R^2=-2$, $RL=0$。任意の三次曲面上の直線 $l$ について
$(l-L/3)^2=-4/3$ なのでHodge指数で $|Rl|\le\sqrt{8/3}<2$。
従って $Pl\ge0$、滑らかな三次del Pezzo曲面の有効曲線錐は直線で生成されるので $P$ はnefである。
$P+L$ は豊富で消滅定理とRiemann--Rochから $h^0(P)=3$。
$|P|$ は平面への双有理射を与える。実際移動部 $U$ では $0<PU\le1$、$0\le U^2\le1$。
$U^2=0$ なら自由で鉛筆の合成となり、二次元線形系なので $U$ は一般繊維の二倍以上となって $PU=1$ に反する
（有理曲面からの底曲線は有理）。$U^2=1$ はHodgeの等号で固定部なしである。
残る一点基点があれば一度吹いた自由系が再び鉛筆と合成され、全引戻し基底で例外係数が $-1$ であることに反する。
従って $P$ は平面線束の引戻しである。縮約成分は負の自乗を持つ三次曲面上の曲線ゆえ随伴により直線で、
互いに素である（負定値性）。この射は六点のブローアップの逆射である。

入力射の正次元繊維の個数を抑える。単純二基点を通る平面束（近接なら指定接方向に沿う平面束）について、
基底の空間直線自体以外の整曲線 $T$ では、$LT$ が両基点での重複度和以上である。
従って $H$ の縮約には $PT=0$ が必要で、これら互いに素な直線のうち基点を通るものは高々二本。
基底の空間直線が曲面にあればそれがさらに一本である。近接の場合に最初の例外直線の固有変換を加えても
入力での例外点は四つ以下である。二重基点なら $H=2(L-e)$ で、縮約はその点からの射影で潰れる直線だけである。
これらは接平面の三次切断に入り、高々三本である。
\end{proof}

\begin{proof}[三点繊維の交点数]
従って長さ三の孤立繊維はそれぞれ少なくとも十六個、十七個である。
一方、平面への表示を使い全引戻し例外類で書くと、二つの入力束は順に
\[
           H=4h-\sum_{i=1}^8 e_i,\qquad
           H=6h-2\sum_{i=1}^7 e_i .
\]
順序付きの中心を一般位置に動かせる。中央の各系は一般切断 $B$ の種数三、$H|_B$ の次数八と
切断完全列より $H^1(S,H)=0$ なので四切断も一緒に変形できる。

一般位置ではどちらの $H$ も非常に豊富である。後者は次数二のdel Pezzo曲面の $-2K_S$ で、
反標準二重平面の二次座標と分岐四次式の平方根の座標による埋込みを与える。
前者は次数一のdel Pezzo上の $h-K_S$ である。Reiderの非常豊富性定理を
$H=K_S+B_0$, $B_0=h-2K_S$ に用いる。$B_0$ は豊富で $B_0^2=17$、
全有効曲線で次数二以上である。残る障害 $B_0T=2,T^2=0$ なら $(-K_S)T=1$ であり、
随伴の偶奇に反する。従って非常に豊富である。
一般の四空間への射影で三点退化軌跡は有限である。実際一般超平面切断は種数三の滑らかな曲線で
入力束の次数八は $2g+2$ 以上、従って完全系の曲線は射影正規でイデアルは二次生成である
（高次数曲線の性質 $N_1$）。$H^1(S,tH)=0$ は $t\ge0$ で切断完全列により帰納的に従い、
曲線への全切断の制限が全射する。よって曲線の射影正規性を超平面列で持ち上げれば曲面も射影正規、
曲線イデアルの二次生成を同じ列で持ち上げれば曲面イデアルも二次生成である。
前定理と同じ平面を張る長さ三スキームの余次元六の議論が使える。

孤立点は一般繊維へ移り、一般の長さは前定理と同じ同伴束の積分である。
平面にトーラスを作用させ、固定点を順次吹くと順序付きブローアップの変形不変性によって
上の二つの束の積分は順に十三、十四である。\path{computations/verify_cofactor_eighteen_cubic.py}
は平面の三つの接表現から出発してこの和を有理数で計算する。いずれも孤立点の下界に反する。
\end{proof}
